\section{Background}
\label{sec:background}

\subsection{Reverso}

Reverso is a zero-shot forecasting foundation model: trained once on a large
corpus of real and synthetic series, then evaluated without fine-tuning. Its
backbone is a DeltaNet-style linear-attention stack
\citep{yang2024deltanet, yang2025gateddeltanet}. Forecasting proceeds by
patches: the model emits $48$ points at a time and the prediction is fed back
autoregressively until the requested horizon is covered. Inputs are min--max
normalised per window, and the output projection therefore lives in normalised
space --- a detail that matters in \S\ref{sec:method:quantile}.

We work with the \texttt{nano} configuration. All comparisons in this report are
\emph{internal}: a baseline trained by us with one variable changed, scored by
our own harness. Two separate reasons make that discipline necessary.

First, our training budget is not theirs. The paper reports compute as
$\{10, 20, 40\}$ H100-hours for its three sizes, so \texttt{nano} is
$10$ H100-hours; our $k{=}1$ arm is about four A100-40GB-hours, i.e.\ roughly
$1.5$--$2$ H100-equivalent hours, so we train \texttt{nano} at
\textbf{15--20\,\% of the paper's compute}. Our reproduction scores
\ReproGap\,\% behind their published checkpoint in our harness
(Table~\ref{tab:arms}), and closing that gap is not what this project is about.

Second, \textbf{our absolute MASE is not on the same scale as the paper's},
and this we have measured. The paper reports \texttt{Reverso-Nano} at $0.757$
on full GIFT-Eval; our harness scores \emph{their own published checkpoint} at
$1.0742$. Dividing our $97$ per-row values by GIFT-Eval's seasonal-naive
scores, on the $97$ rows both sides share, gives

\begin{center}\small
\begin{tabular}{lrr}
  \toprule
   & raw & normalised \\
  \midrule
  all 97 rows & 1.0742 & 0.7684 \\
  short (55)  & 1.1204 & 0.7594 \\
  medium (21) & 0.9988 & 0.7731 \\
  long (21)   & 1.0346 & 0.7877 \\
  \midrule
  seasonal-naive baseline & 1.3979 & --- \\
  \bottomrule
\end{tabular}
\end{center}

So the factor is the benchmark's aggregation convention: GIFT-Eval's
leaderboard divides each row's MASE by the seasonal-naive baseline's before
taking the geometric mean, while we report raw MASE. Raw MASE exceeds $1$ for
every model because its denominator is a one-step \emph{in-sample} error while
its numerator is a $48$-step \emph{out-of-sample} one. Nothing about our
training is implicated.

This matters for presentation only, never for a comparison, because a per-row
divisor cancels in a ratio of geometric means: with
$\mathrm{MASE}_{\text{norm}} = \mathrm{MASE}/S$ we have
$\operatorname{geo}(A/S) = \operatorname{geo}(A)/\operatorname{geo}(S)$, so
\[
  \frac{\operatorname{geo}(A/S)}{\operatorname{geo}(B/S)}
  \;=\; \frac{\operatorname{geo}(A)}{\operatorname{geo}(B)} .
\]
Every percentage in this report is therefore invariant to the convention. We
report raw MASE throughout, and print no absolute figure beside a published
one.

\subsection{A residual we cannot yet explain, and where it probably comes from}
\label{sec:background:backend}

The normalisation does not close the gap entirely. Normalised, we score the
authors' checkpoint $1.51\,\%$ worse than they publish ($0.7684$ against
$0.757$), and $2.84\,\%$ worse on the long horizon ($0.7877$ against their
$0.766$) --- a discrepancy that \emph{grows with horizon}. That is larger than
the looping effect this report measures, so it deserves naming rather than
rounding away.

It has a concrete candidate. We have never evaluated with the authors'
implementation: \texttt{eval\_gift.py} runs our own \texttt{reverso\_torch}
re-implementation, because the official package imports FlashFFTConv at module
load and the cluster's \texttt{nvcc}~11.7 does not match torch's CUDA~12.4. The
weights load soundly --- only the \texttt{flashfftconv.*} buffers are stripped,
then \texttt{strict=True}, exactly as upstream --- so nothing is silently
dropped. What differs is the long convolution. Theirs is
\texttt{FlashFFTConv(seq\_len)} with the pre-gate fused into the kernel; ours is
an FFT product at length $N$, that is a \emph{circular} convolution, carrying a
comment asserting it matches the kernel. The long convolutions FlashFFTConv was
designed for are \emph{causal}: zero-pad to $2N$, keep the first $N$. If that is
the kernel's semantics, our version wraps the oldest part of the context onto
the most recent positions --- precisely where a forecaster looks hardest, and
precisely the error that compounds along an autoregressive rollout, which is
what a horizon-growing discrepancy looks like.

It is also not a constant offset, because the loop iterates the whole layer
stack, convolution included: at $k{=}4$ the artifact is applied four times. So
it is a candidate confound for \S\ref{sec:results}, not merely a calibration
issue.

\todo{next experiment, before the final report}{%
  Re-evaluate the \emph{published} checkpoint with the zero-padded
  ($n = 2N$, first $N$) convolution. No retraining and no FlashFFTConv needed;
  one evaluation, about two GPU-hours. If the normalised score moves toward
  $0.757$ and the long term toward $0.766$, the comment is wrong and the
  finding belongs in \S\ref{sec:discussion} as a limitation on what model these
  results describe. Our own arms must \emph{not} be re-scored under the new
  convolution --- they were trained with the circular one
  (\S\ref{sec:protocol}).}

\subsection{The delta rule and its spectrum}

DeltaNet maintains a matrix state $S_t$ updated by the delta rule
\begin{equation}
  S_t \;=\; \bigl(I - \beta_t\, k_t k_t^{\!\top}\bigr) S_{t-1}
          \;+\; \beta_t\, k_t v_t^{\!\top},
  \label{eq:deltarule}
\end{equation}
with $k_t$ $L_2$-normalised. The transition matrix
$I - \beta_t k_t k_t^{\!\top}$ is a rank-one perturbation of the identity: its
eigenvalues are $1$ on the subspace orthogonal to $k_t$ and $1-\beta_t$ along
$k_t$. With the usual $\beta_t \in (0,1)$ the state is \emph{contracted} along
the key direction and never reflected, so the reachable transitions form a
strictly contractive family.

This spectral reading is what makes one of our three levers interesting rather
than arbitrary. Permitting $\beta_t \in (0,2)$ makes $1-\beta_t$ negative: the
state is \emph{reflected} along $k_t$ instead of merely shrunk, while
$|1-\beta_t|\le 1$ keeps the map non-expansive, so training stability is
preserved. \citet{grazzi2025unlocking} showed that this extended range strictly
enlarges the class of functions a linear-attention layer can express, and
\citet{sarrof2024expressive} that a signed gate is needed to represent
alternating structure at all.

Very recent work makes the geometric point sharp. A rotation in a plane is the
composition of \emph{two} reflections, and a 2D rotation is exactly what a
recurrence needs in order to carry oscillatory structure forward. One can get
such a rotation by composing two delta-rule transitions in a single update, but
that doubles the rank and the cost. Complex KDA \citep{ckda2026} shows that
Kimi Delta Attention already contains the second reflection: the channel-wise
gate supplies it, provided the gate is allowed to be \emph{negative}. Combining
the extended $\beta \in [0,2]$ of \citet{grazzi2025unlocking} with a signed
channel-wise decay $\alpha \in [-1,1]^n$ gives the transition
\begin{equation}
  A_t \;=\; \bigl(I - \beta_t k_t k_t^{\!\top}\bigr)\,\mathrm{Diag}(\alpha),
  \label{eq:ckda}
\end{equation}
whose spectrum can leave the real line at no extra rank. Their Table~1 states
the boundary exactly: the number of complex eigenvalue pairs is \emph{zero}
whenever $\beta \le 1$ or $\alpha \ge 0$, and a \emph{scalar} gate cannot help,
because scalars commute with everything and so cannot contribute a reflection
in a chosen direction. The reported payoff is length extrapolation on $S_3$,
$S_4$ and \emph{periodic waveform continuation} --- which is, almost literally,
seasonal forecasting. Two of that paper's authors supervise or lead this lab,
and the implementation is a minor modification of the FLA kernels we already
run, at $96$--$97\,\%$ throughput.

Our \texttt{allow\_neg\_eigval} lever (\S\ref{sec:method:levers}) is the
\emph{first} of those two ingredients only: it supplies one reflection and
leaves the decay unsigned, so by \eqref{eq:ckda} its spectrum stays real. That
it already buys $+2.22\,\%$, with the only horizon-growing profile among our
three levers, is what makes the second ingredient the most promising remaining
experiment (\S\ref{sec:remaining}).

\subsection{Looped and weight-tied models}

Weight tying across depth goes back to the Universal Transformer
\citep{dehghani2019universal}. The recent revival comes from reasoning: looped
transformers can emulate iterative algorithms \citep{giannou2023looped}, and
looped language models trained with iterative latent computation match much
larger models on reasoning benchmarks \citep{zhu2026ouro}. The claim in that
literature is explicitly about \emph{latent reasoning} --- more iterations,
more steps of an algorithm --- which is the part that does not transfer
obviously to forecasting, and which this project tests.

Closest to our setting is \citet{moroshan2026tempopfn}, which pretrains linear
RNNs on synthetic data for zero-shot forecasting: same model family, same
benchmark, no looping.

\subsection{GIFT-Eval}

GIFT-Eval \citep{aksu2024giftevalbenchmark} is a zero-shot forecasting
benchmark spanning many domains and three horizon regimes. We report the
geometric mean of MASE over the $97$ (dataset, horizon) rows the model
supports, split $55$ short / $21$ medium / $21$ long. The geometric mean is the
benchmark's own convention and it matters here: it keeps the aggregate robust
to a single pathological dataset, which a looped model evaluated at the wrong
$k$ produces easily (\S\ref{sec:protocol}).
